% !TEX TS-program = pdflatexmk
\documentclass{article}
\usepackage[english]{babel}
\usepackage[letterpaper,top=2cm,bottom=2cm,left=3cm,right=3cm,marginparwidth=1.75cm]{geometry}
\usepackage{csquotes}
\usepackage[
  backend=biber,
  style=apa,
  sortcites=true,
  natbib=true,
  maxcitenames=2,   % APA-like in-text truncation
  maxbibnames=99
]{biblatex}
\DeclareLanguageMapping{english}{english-apa}
\addbibresource{references.bib}
\usepackage{amsmath}
\usepackage{float}
\usepackage{graphicx}
\usepackage{siunitx}
\usepackage[colorlinks=true, allcolors=blue]{hyperref}
\usepackage{booktabs}
\usepackage{caption}
\usepackage{dcolumn}          % decimal-aligned columns for regression tables
\usepackage{threeparttable}   % table notes

\title{The Price Elasticity of Demand for Marijuana in Mexico}
\author{Alejandro Romero Gonz\'{a}lez}

\begin{document}
\maketitle

\begin{abstract}
Mexico's cannabis-policy debate requires evidence on how consumers respond to prices, yet transaction-level price and quantity data are rarely available in illicit markets.  This paper estimates the transaction-level price elasticity of marijuana demand in Mexico using an anonymous crowdsourced survey of 3{,}293 purchases recorded in 2019 across all 32 states and 351 municipalities.  In log--log specifications with product-quality controls, demographics, and state-by-month fixed effects, the benchmark elasticity is $-0.652$: a 10\% higher unit price is associated with a 6.5\% smaller purchase quantity.  Demand differs sharply by quality tier, with seeded marijuana more price responsive ($-0.891$) than seedless marijuana ($-0.613$).  The estimate remains similar under extreme-value trimming, heaping checks, municipality fixed effects, wild cluster bootstrap inference, and coefficient-stability bounds.  Instrumental-variable diagnostics qualify rather than replace the benchmark: OSM production-hub driving time remains weak, hub-only lagged seizure gravity reproduces the benchmark magnitude, and external Hausman instruments that exclude the buyer's own state or macroregion have strong first stages and yield negative elasticities.  The results provide the first Mexico-specific transaction-level benchmark for cannabis tax and regulatory analysis, while emphasizing that illicit-market elasticities should not be mechanically extrapolated to participation, aggregate consumption, or a future regulated market.
\end{abstract}

%% ---- Sections ----
\section{Introduction}\label{sec:intro}
% ============================================================
%  sections/intro.tex  –  Introduction
% ============================================================

Over the last two decades, Mexico has repeatedly revisited cannabis policy through criminal-law reform, public-health legislation, and constitutional adjudication.  The 2009 \emph{narcomenudeo} reform decriminalized small-quantity possession and reassigned retail drug offenses to state courts \citep{dof2009narcomenudeo,lgs2009doseTable}.  In 2017, amendments to the General Health Law authorized medical and scientific uses \citep{dof2017medicalcannabis}.  Most recently, the Supreme Court's 2021 General Declaration of Unconstitutionality (DGI~1/2018) invalidated the blanket prohibition on recreational use, mandating comprehensive regulation \citep{scjn2021dgi}.

These legal shifts have been accompanied by sustained attention to the fiscal design of a regulated market.  The Senate approved in 2020 the \emph{Ley Federal para la Regulaci\'{o}n del Cannabis}, whose fiscal chapter, building on an earlier legislative initiative \citep{delgado2019iniciativa}, envisioned a 12\% \emph{ad valorem} excise tax plus a per-gram surcharge under the IEPS framework \citep{senado2021minuta}.  The Chamber of Deputies passed an amended version in March~2021, but disagreements between the two chambers left the bill without definitive enactment \citep{diputados2021dictamen}.  Comparative analyses have evaluated international cannabis-tax regimes to inform domestic design \citep{clavellina2020cannabis_taxes}, academic work has projected the fiscal-revenue implications of legalization \citep{hernandezmacías2021impacto}, and microeconomic studies have examined the theoretical conditions for optimal taxation of marijuana in Mexico \citep{perezromero2014elementos}.  Designing any such tax regime, and more broadly anticipating the market consequences of regulation, requires an empirical input that, to date, does not exist for Mexico: the price elasticity of demand for marijuana.

The absence of domestic estimates reflects a fundamental data constraint.  In illicit markets, transaction-level prices and quantities are rarely observed at scale; administrative price series are either absent or unreliable \citep{caulkins1998price}, and household surveys such as Mexico's ENCODAT capture consumption prevalence but not the price--quantity margin \citep{encodat2025}.  The international literature offers meta-analytic consensus elasticities in the range of $-0.30$ to $-0.56$ \citep{gallet2014monkey}, complemented by crowdsourced estimates for the United States \citep{davis2016crowd}, regionally disaggregated analyses \citep{halcoussis2017geo}, online-survey evidence from South Africa \citep{riley2020southafrica}, and recent evidence from legal retail markets \citep{han2025commercialization,dodd2025canada}.  What remains missing is a Mexico-specific estimate based on realized transactions, in a market where policy reform, supply-chain geography, quality segmentation, and enforcement conditions differ from the settings studied to date.

This paper fills the gap by exploiting a rare data source: an anonymous survey, fielded between February~19 and September~15, 2019 and promoted mainly through \emph{Vice en Espa\~{n}ol}, that asked respondents to report completed marijuana purchases in Mexico \citep{bejaranoromero2020dataset}.  The raw file includes response identifiers and timestamps; respondent age, gender, education, and self-reported socioeconomic status; state and municipality of purchase; purchase date; perceived product quality (seedless or seeded); quality-specific quantity and expenditure fields; the buyer's relationship to the seller; age of first cannabis consumption; purchase frequency; and usual purchase amount.  For each retained record, we coalesce the seeded and seedless quantity/expenditure fields into one transaction-level quantity and value, then compute unit price as pesos per gram.  After cleaning, the analytic sample contains 3{,}293 purchases dated January~1--September~11, 2019, spanning all 32 states and 351 municipalities.  The survey is not nationally representative; its value is that it observes realized prices, quantities, quality, timing, and geography at a level of detail unavailable in conventional Mexican data sources.

We estimate log--log demand specifications that yield a stable transaction-level price elasticity of approximately $-0.65$ across OLS models with progressively richer controls. This elasticity describes how the quantity purchased in a completed transaction varies with the unit price recorded for that transaction; it does not directly identify participation responses, purchase frequency, or total individual consumption.  Quality segmentation reveals meaningful heterogeneity: seedless (\emph{good}) marijuana has an elasticity of about $-0.61$, whereas seeded (\emph{bad}) marijuana is more price responsive at roughly $-0.89$.  The benchmark estimate remains similar under extreme-value trimming, heaping checks, municipality fixed effects, wild cluster bootstrap inference, and coefficient-stability diagnostics.  We also implement instrumental-variables diagnostics in three families: OpenStreetMap (OSM) road-network driving time to major production states computed with the Open Source Routing Machine (OSRM); one-month-lag seizure-gravity measures based on OSM driving time and M\'{e}xico Unido contra la Delincuencia (MUCD) marijuana-seizure data, derived from official responses to public-information requests; and external Hausman price instruments that average prices in other markets within the same quality--month cell.  The OSM production-hub instrument remains too weak to serve as an alternative benchmark.  The hub-only lagged seizure-gravity specification reproduces the benchmark magnitude, while stricter outside-region seizure gravity illustrates the tradeoff between a cleaner exclusion story and weaker first-stage strength.  The external Hausman instruments are stronger and produce negative elasticities.  We interpret these IV results as diagnostics that triangulate the benchmark rather than as a definitive causal replacement for it.

The study contributes to the literature in three ways.  First, it provides the first Mexico-specific transaction-level estimate of the price elasticity of marijuana demand, directly informing a policy debate in which cannabis taxation and regulation have advanced faster than domestic demand evidence \citep{perezromero2014elementos,clavellina2020cannabis_taxes,senado2021minuta}.  Second, it shows that quality-tier heterogeneity is empirically important in this setting: seeded and seedless marijuana differ not only in prices but also in estimated demand responsiveness.  Third, it evaluates the promise and limits of crowdsourced transaction data for illicit-market demand estimation by pairing the survey with municipality-level geography, OSM/OSRM driving times, and MUCD enforcement records \citep{davis2016crowd,halcoussis2017geo,riley2020southafrica}.  The policy relevance is therefore practical but bounded.  An elasticity near $-0.65$ suggests that price-based regulation can affect purchase quantities without eliminating the tax base, while the stronger response for seeded marijuana warns that uniform taxes may affect quality tiers unevenly.  At the same time, the estimate comes from illicit-market transactions before national recreational regulation, so it should be treated as one empirical input into tax design rather than as a complete forecast of a legalized market.

The remainder of the paper is organized as follows: Section~\ref{sec:litreview} reviews the literature, Section~\ref{sec:data} describes the data and empirical strategy, Section~\ref{sec:results} presents the results, and Section~\ref{sec:conclusions} concludes with policy implications.


\section{Literature Review}\label{sec:litreview}
% ============================================================
%  sections/litreview.tex  –  Literature Review
% ============================================================

The international evidence consistently indicates a negative price elasticity of demand for cannabis, but estimates span a wide range depending on the behavioral margin examined, the measurement strategy, and the institutional setting in which transactions occur.  This section organizes the relevant literature along these dimensions (meta-analytic syntheses, micro-level measurement, crowdsourced transaction data, cross-country evidence with quality controls, the participation margin, legalization-era identification, and Mexico-specific analysis) and concludes by motivating the contribution of the present paper.

The broadest synthesis is the meta-analysis of illicit drug demand by \citet{gallet2014monkey}, which compiles hundreds of elasticity estimates across substances and reports a median price elasticity of demand for marijuana of $-0.33$.  Meta-regressions in that study reveal systematic heterogeneity: long-run elasticities tend to exceed short-run elasticities in absolute value, participation-margin estimates are generally more negative than intensive-margin estimates, and studies relying on hypothetical price scenarios produce substantially larger responses than those using observed transactions.  These patterns underscore the value of transaction-level data and motivate caution when extrapolating across settings.

Early micro evidence illustrates how the choice of data source shapes estimated responsiveness even within a single population.  Using survey data on students at the University of California, Los Angeles, \citet{nisbet1972ucla} estimate elasticities from both a hypothetical market-survey demand schedule and realized purchase records.  The double-log specification yields an elasticity of $-1.013$ for realized purchases versus $-0.365$ for the market survey; evaluated at the then-prevailing price of \$10 in a linear model, the analogous figures are $-1.51$ and $-0.51$.  The gap is directly relevant for illicit-market analysis: stated demand schedules may fail to capture search costs, legal risk, and other frictions, whereas realized transactions incorporate the constraints buyers actually face.  This early finding foreshadows a persistent methodological lesson in the literature, namely that the data-generating process matters as much as the econometric specification for recovering behaviorally meaningful demand parameters.

More recent work exploits transaction-level information from crowdsourced platforms and typically finds inelastic but economically meaningful responses.  \citet{davis2016crowd} use crowdsourced U.S.\ transaction data from \texttt{PriceOfWeed.com}, an online price-reporting platform, and estimate an OLS elasticity of $-0.689$.  To address potential price endogeneity, they employ instrumental-variables strategies using plausibly supply-driven shifters (electricity prices and geographic distance to major growing regions) and report IV estimates ranging from $-0.597$ to $-0.789$ depending on the instrument set.  The identification logic is that these cost shifters affect the retail price through the supply side but, conditional on controls, do not independently influence the quantity a given buyer demands.  Similarly, \citet{halcoussis2017geo} use \texttt{PriceOfWeed.com} data to model price and quantity jointly using geographic cost variation tied to distance from production areas, reporting a price elasticity of $-0.418$.  Extending the crowdsourced approach to a different national context, \citet{giommoni2018uk} analyze the British cannabis market using \texttt{PriceOfWeed.com} reports and document substantial regional price variation, providing further evidence that crowdsourced transaction data can generate economically informative estimates of illicit-market structure.  Together, these studies suggest that transaction-based elasticities for cannabis typically fall in the range of roughly $-0.4$ to $-0.8$, while highlighting that credibility depends on the defensibility of exclusion restrictions.

Survey-based transaction evidence from outside the United States points to comparable magnitudes and underscores the importance of controlling for product quality.  Using an online survey of South African cannabis users, \citet{riley2020southafrica} estimate a baseline conditional elasticity of $-0.501$.  Across alternative specifications that vary the controls for product differentiation, the coefficient ranges from $-0.493$ to $-0.613$.  This sensitivity to quality-related controls implies that, in markets where product heterogeneity is pronounced, omitting quality proxies risks biasing the estimated demand response.  The concern is directly relevant here because, as shown below, seedless and seeded marijuana command sharply different prices in the Mexican survey data.  \citet{benlakhdar2016potency} reinforce this point using data from over 250 regular cannabis users in France, estimating short-run price elasticities between $-1.7$ and $-2.1$.  Notably, controlling for real or perceived potency has little effect on the magnitude of the price response in their sample, suggesting that while quality matters for willingness to pay, the price elasticity itself may be relatively robust to potency measurement once product heterogeneity is accounted for.

A related strand of the literature estimates elasticities on the extensive (participation) margin.  \citet{pacula2000youth} show that youth participation elasticities are sensitive to trend specification and the inclusion of mediators, reinforcing the distinction between participation and intensive-margin responses.  \citet{amlung2019substitution} complement this work by documenting that demand for illegal cannabis is relatively more elastic than demand for legal cannabis, implying that consumers are willing to switch sources as relative prices shift.

Recent legalization in several U.S.\ states and in Canada has created new opportunities for identification via fiscal-policy and regulatory variation.  \citet{han2025commercialization} study adult past-30-day use in states with recreational commercialization and exploit tax rates as an instrument for retail prices, reporting implied elasticities between $-0.66$ and $-0.59$.  \citet{dodd2025canada} provide the first post-legalization elasticity estimate for Canada's regulated market, using sales data from British Columbia and reporting a price elasticity of approximately $-1.4$, substantially more elastic than most estimates from illicit settings.  Their approach illustrates the appeal of regulated-market data, where prices, quantities, and product attributes are reliably observed, but the resulting estimates may not transfer directly to illicit markets where administered tax variation is absent, product quality is heterogeneous, and buyers face additional legal and search-cost frictions.

Mexico-specific work has engaged with the taxation question but largely without transaction-level evidence capable of identifying a domestic price elasticity.  \citet{perezromero2014elementos} develop a microeconomic framework for analyzing cannabis legalization and indirect taxation.  A central insight of their analysis is that the relevant price faced by consumers in an illicit market should be understood as a ``full price'' encompassing not only the monetary cost but also legal risk, search costs, and other transaction frictions associated with purchasing from an illegal supplier.  They further discuss how substitution and complementarity with legal substances, notably alcohol and tobacco, may shape the effects of legalization on consumption patterns.  At the same time, their discussion highlights a binding constraint: in the absence of reliable price and consumption data from the Mexican market, direct estimation of demand parameters has not been feasible.

Two points follow from this literature.  First, cannabis demand slopes downward and is often less than unit-elastic, but estimates vary across behavioral margins, measurement strategies, and institutional settings \citep{gallet2014monkey}.  Second, credible estimation requires explicit attention to price endogeneity and product differentiation \citep{davis2016crowd,halcoussis2017geo,riley2020southafrica,benlakhdar2016potency}.  No study to date has estimated the transaction-level price elasticity of marijuana demand using purchase records from Mexico.  This paper fills that gap with a raw crowdsourced survey of 4{,}775 purchase records that jointly observe price, quantity, perceived quality, and fine-grained geography \citep{bejaranoromero2020dataset}; after cleaning, the analytic sample retains 3{,}293 transactions.  The log--log demand framework controls for quality tiers, captures the sharp price wedge between seedless and seeded product documented in the data, and absorbs unobserved geography and timing through fixed effects, with the preferred specification using state-by-month fixed effects.

Price endogeneity is a first-order concern in any demand analysis, and illicit markets pose additional identification challenges because standard supply-side instruments (taxes, regulated input costs) are unavailable \citep{han2025commercialization,dodd2025canada}.  We address this concern along two complementary dimensions.  First, our baseline specifications systematically vary controls, fixed-effects structures, and sample restrictions, allowing us to assess coefficient stability and the sensitivity of the estimated elasticity to potential confounders; we formalize this approach via the \citet{oster2019unobservable} bounding procedure.  Second, we implement instrumental-variables strategies motivated by the geography-based cost shifters used by \citet{davis2016crowd} and \citet{halcoussis2017geo}, constructing instruments based on Open Source Routing Machine (OSRM) driving times over OSM road data to major production hubs, a one-month-lag OSM driving-time-weighted seizure index drawn from M\'{e}xico Unido contra la Delincuencia (MUCD) municipality-month enforcement data \citep{openstreetmap2026,luxen2011osrm,mucd2019seizures}, and a Hausman-type leave-one-out mean price.  We report these IV results primarily as diagnostics: the Hausman-type instrument provides a strong internal robustness check under its leave-one-out exclusion restriction, while the external geography-based instruments require weak-IV diagnostics because their first stages are limited.  This transparent, assumption-explicit approach provides an empirical foundation for the ongoing policy discussion on cannabis regulation and taxation in Mexico.


\section{Data, Descriptive Statistics, and Empirical Strategy}\label{sec:data}
% ============================================================
%  sections/data.tex  –  Data, Descriptive Statistics, and Empirical Strategy
% ============================================================

\subsection{Data Source}

The analysis presented in this paper draws on a crowdsourced dataset of marijuana purchases in Mexico assembled by \citet{bejaranoromero2020dataset}.  An anonymous online questionnaire, administered via Qualtrics from February~19 to September~15, 2019 and mainly promoted through the news outlet \emph{Vice en Espa\~{n}ol}, collected self-reported information on individual transactions.  For each purchase, the respondent recorded the quantity acquired, the amount paid, perceived product quality, the state and municipality of purchase, and the purchase date.  The questionnaire also asked about self-reported socioeconomic status, age of first cannabis consumption, purchase frequency, usual purchase amount, the respondent's relationship with the dealer, and basic demographic characteristics.  The raw dataset contains 4{,}775~records.

We supplement the survey with two external spatial and enforcement data sources.  Municipality and state identifiers are linked to INEGI municipal and state polygons from the Marco Geoestad\'{i}stico Integrado; we calculate municipal centroids for purchase and seizure municipalities and full-state centroids for the production-hub states used in the OSM driving-time instrument \citep{inegi2022mgi}.  Road-network driving times are computed with the Open Source Routing Machine (OSRM) using OpenStreetMap (OSM) data \citep{openstreetmap2026,luxen2011osrm}.  Because the OSRM matrices are built from an OSM road-network snapshot queried in April~2026, these variables should be interpreted as persistent road-network connectivity rather than contemporaneous 2019 traffic conditions.  Enforcement intensity is measured with monthly municipality-level data published by M\'{e}xico Unido contra la Delincuencia (MUCD), a Mexican non-governmental organization.  These records report marijuana seizures, crop destruction, and seed confiscation; the seizure-gravity instrument uses only marijuana seizures from the one-month lag window relevant for the estimation sample, December~2018 through August~2019 \citep{mucd2019seizures}.  MUCD's platform compiles official responses obtained through public-information requests to the transparency units of the Secretariat of the Navy (SEMAR), the Secretariat of National Defense (SEDENA), the National Guard (GN), the former Federal Police (PF), and the Attorney General's Office (FGR); MUCD then transcribed, systematized, cleaned, and fused the heterogeneous source files into a municipality-month database \citep{mucd2019seizures}.  These external sources are used only to construct the instrumental-variables diagnostics; the baseline OLS estimates rely on the cleaned transaction-level survey sample.  Appendix~\ref{app:external-data} reports descriptive diagnostics for the MUCD marijuana-seizure records used in the IV diagnostics and the OSM/OSRM driving-time matrices.

These auxiliary sources characterize geography and enforcement, but the core empirical advantage comes from the survey's transaction-level design.  Conventional data sources usually miss the relevant price--quantity margin: administrative price series do not exist for illegal drugs, and household surveys such as ENCODAT measure consumption prevalence but generally do not record the price and quantity of specific purchases \citep{encodat2025}.  Here, respondents report completed transactions, so the data reflect the terms buyers actually faced rather than hypothetical willingness to pay.  That distinction matters for demand estimation: search costs, legal risk, and imperfect information about quality can all affect realized purchases, and \citet{nisbet1972ucla} show that elasticities based on actual purchases can differ substantially from hypothetical-response estimates.

\subsection{Sample Construction}

The analytic sample is built around the variables needed to observe a completed purchase and estimate a revealed-demand elasticity.  Quantity purchased is the outcome, while transaction value and quantity jointly define the unit price faced by the buyer.  Perceived quality is retained because seeded and seedless marijuana are not the same product, and failing to condition on quality would confound price variation with product differentiation.  Purchase date and municipality/state identifiers are required to align transactions with local market conditions, fixed effects, and the external driving-time and enforcement instruments.  Finally, age, gender, and education are included as baseline demographic controls because they are predetermined respondent characteristics that can shape purchasing behavior, income constraints, risk exposure, and access to suppliers.

The survey also records self-reported socioeconomic status, age of first consumption, purchase frequency, usual purchase amounts, and the respondent's relationship with the dealer.  These variables are informative about consumer experience, market relationships, and socioeconomic position, but we do not use them as baseline controls or sample-defining fields.  Frequency, usual amount, and buyer--seller relationship are closer to behavioral intensity or matching outcomes than to predetermined covariates, and conditioning on them could absorb part of the demand response or the negotiation channel through which prices and quantities are jointly determined.  Socioeconomic status is self-classified and coarser than education, so we treat it as contextual rather than as a main control.  Requiring complete responses for these auxiliary items would also reduce the sample without being necessary for the core price--quantity estimand.

We apply a sequence of filters and transformations to arrive at the analytic sample. First, we require nonmissing values in the fields needed by the main specification: state, municipality, age, gender, education, purchase date, perceived quality, quantity, and transaction value. This complete-case requirement removes 1{,}291 of the 4{,}775 raw records. Second, we exclude 7 additional records with implausible age, transaction-value, or quantity entries. Third, although the questionnaire was fielded from February to September~2019, some respondents reported earlier purchases; restricting to transactions dated on or after January~1, 2019 removes 184 recalled purchases. We then construct the unit price per gram as the ratio of transaction value to quantity purchased, $p_i = \text{value}_i / \text{quantity}_i$, and code perceived quality as a binary indicator, with \emph{good} denoting seedless marijuana and \emph{bad} denoting seeded marijuana. After all filters and transformations, the analytic sample comprises 3{,}293 observations, with retained purchases dated January~1--September~11, 2019.

\subsection{Descriptive Statistics}

Table~\ref{tab:descriptive} presents summary statistics for the full sample and by quality tier.  The median transaction involves 14~grams purchased for 250~MXN, yielding a median price of about 21~MXN per gram.  The distributions of both price and quantity are right-skewed, with means substantially exceeding medians and large standard deviations relative to the mean, motivating the log transformation used in the regressions.  The median respondent is~25 years old.

\begin{table}[H]
\centering
\caption{Descriptive statistics}
\label{tab:descriptive}
\begin{threeparttable}
\begin{tabular}{lrrrrr}
\toprule
 & Mean & Median & SD & Min & Max \\
\midrule
\multicolumn{6}{l}{\textit{Panel A: Full sample (N = 3,293)}} \\[3pt]
Price (MXN per gram)   & 41.97 & 21.43 & 88.35 & 0.01 & 3,000.00 \\
Quantity (grams)       & 31.27 & 14.00 & 83.40 & 1.00 & 2,000.00 \\
Transaction value (MXN)& 568.03 & 250.00 & 1,892.43 & 1.00 & 65,000.00 \\
Age (years)            & 26.20 & 25.00 & 6.74 & 18 & 68 \\
\\[6pt]
\multicolumn{6}{l}{\textit{Panel B: Good quality (seedless) (N = 2,636)}} \\[3pt]
Price (MXN per gram)   & 48.39 & 28.57 & 97.12 & 0.01 & 3,000.00 \\
Quantity (grams)       & 31.31 & 14.00 & 83.54 & 1.00 & 2,000.00 \\
Transaction value (MXN)& 663.33 & 320.00 & 2,100.50 & 1.00 & 65,000.00 \\
Age (years)            & 26.38 & 25.00 & 6.75 & 18 & 65 \\
\\[6pt]
\multicolumn{6}{l}{\textit{Panel C: Bad quality (seeded) (N = 657)}} \\[3pt]
Price (MXN per gram)   & 16.24 & 10.00 & 21.29 & 0.25 & 200.00 \\
Quantity (grams)       & 31.13 & 10.00 & 82.90 & 1.00 & 1,000.00 \\
Transaction value (MXN)& 185.68 & 100.00 & 258.32 & 10.00 & 2,500.00 \\
Age (years)            & 25.47 & 24.00 & 6.68 & 18 & 68 \\
\bottomrule
\end{tabular}
\begin{tablenotes}[flushleft]\small
\item \textit{Notes:} Price per gram is computed as transaction value divided by quantity purchased.  Quality classification is self-reported: ``good'' denotes seedless marijuana and ``bad'' denotes seeded marijuana.
\end{tablenotes}
\end{threeparttable}
\end{table}


Quality segmentation is pronounced.  Seedless (\emph{good}) marijuana, which accounts for 80\% of transactions ($N = 2{,}636$), commands a median price nearly three times that of seeded (\emph{bad}) marijuana (28.57 versus 10.00~MXN per gram).  The price distribution for bad-quality product is considerably more compressed (SD~=~21 versus~97), consistent with a less differentiated, commodity-like segment.  Median quantities are also higher for good-quality purchases (14 versus 10~grams), though the means are similar across tiers.

The sample is predominantly male (84\%), with women comprising 14\% and other gender identities 2\%.  Educational attainment skews toward higher levels: 51\% report a bachelor's degree, 37\% high school, 9\% a graduate degree, and only 3\% middle school or below.  These demographics are consistent with the online, media-promoted recruitment channel and should be borne in mind when considering external validity.

Geographically, the data cover all 32~states and 351~municipalities.  Mexico City dominates with 1{,}011~observations (31\%), followed by the State of Mexico (283), Jalisco (226), and Nuevo Le\'{o}n (150).  The temporal distribution peaks in February and March~2019 (1{,}089 and 903~observations, respectively), coinciding with the survey's launch and initial media push, and tapers in subsequent months.

Quantity heaping is evident: 42\% of reported quantities are multiples of~5, 23\% are multiples of~10, and 14\% are close to~28~grams (approximately one ounce).  This pattern, common in self-reported survey data, motivates a robustness check that excludes round-number quantities.

\subsection{Empirical Strategy}

\paragraph{Baseline specification.}
We estimate the price elasticity of demand using a log--log specification:
\begin{equation}\label{eq:loglog}
  \ln Q_i \;=\; \alpha \;+\; \beta \, \ln P_i \;+\; \gamma \, \textit{Good}_i \;+\; \mathbf{X}_i' \boldsymbol{\delta} \;+\; \mu_s \;+\; \varepsilon_i,
\end{equation}
where $Q_i$ is the quantity purchased (grams), $P_i$ is the unit price (MXN per gram), $\textit{Good}_i$ is an indicator for seedless quality, $\mathbf{X}_i$ is a vector of demographic controls (age, gender, and education), $\mu_s$ denotes fixed effects, and $\varepsilon_i$ is the error term.  The coefficient~$\beta$ is the price elasticity of demand: a one-percent increase in price is associated with a $\beta$-percent change in quantity demanded.

We progressively enrich the fixed-effects structure: (i)~no fixed effects, (ii)~state fixed effects ($\mu_s$), (iii)~state plus month fixed effects, (iv)~state-by-month interactions, and (v)~municipality fixed effects.  The motivation for this progression is to absorb increasingly fine-grained unobserved heterogeneity in local market conditions, enforcement intensity, and demand-side factors.  For the pooled specifications we report heteroskedasticity-robust standard errors; once geographic fixed effects are introduced, standard errors are clustered at the state level.  Because there are 32 state clusters, we also report wild cluster bootstrap inference for the benchmark specification.

\paragraph{Quality heterogeneity.}
To allow the demand response to differ across product tiers, we estimate separate regressions for good and bad quality subsamples, and a pooled specification that interacts $\ln P_i$ with the quality indicator:
\begin{equation}\label{eq:interaction}
  \ln Q_i \;=\; \alpha \;+\; \beta_0 \, \ln P_i \;+\; \beta_1 \, (\ln P_i \times \textit{Good}_i) \;+\; \gamma \, \textit{Good}_i \;+\; \mathbf{X}_i' \boldsymbol{\delta} \;+\; \mu_s \;+\; \varepsilon_i.
\end{equation}
If seedless and seeded marijuana occupy distinct market segments with different demand structures, the interaction coefficient $\beta_1$ will be nonzero.

\paragraph{Identification.}
The key identification challenge is that price may be correlated with unobserved demand-side factors.  In illicit markets, experienced users who purchase frequently may negotiate lower prices, generating a spurious positive correlation between quantity and the inverse of price that biases the OLS elasticity toward zero.  Conversely, buyers in high-demand locations may face elevated prices if local supply is inelastic, biasing the elasticity away from zero.

A second concern is measurement error in self-reported quantities and transaction values, especially because unit price is constructed as transaction value divided by quantity.  Asking value and quantity as separate transaction fields reduces reliance on a single reported unit-price measure, but it does not eliminate mechanical-ratio concerns if quantity is mismeasured.  We therefore report robustness checks that trim extreme price and quantity observations and exclude round-number quantities, the most visible form of heaping in the raw data.

These issues motivate an identification strategy that combines measurement-error checks with approaches to price endogeneity.  We first examine whether the estimated elasticity remains stable across progressively demanding fixed-effects structures, from no controls through state-by-month interactions.  We then implement three instrumental-variables strategies that exploit plausibly supply-driven variation in prices:

\begin{enumerate}
  \item \textbf{OSM driving time from production hubs.}  Following the supply-cost logic of \citet{davis2016crowd} and \citet{halcoussis2017geo}, we instrument price with the log driving time from the nearest major marijuana-producing state (Sinaloa, Chihuahua, Guerrero, Durango, and Michoac\'{a}n) to the purchase municipality.  The selection combines the traditional ``Golden Triangle'' production region of Sinaloa--Durango--Chihuahua with Guerrero and Michoac\'{a}n, which enforcement and cultivation assessments identify among Mexico's main marijuana-growing areas \citep{medel2015cultivation,ndic2008ndta}.  We calculate municipality centroids from projected INEGI municipal polygons and construct each production-hub point as the centroid of the corresponding projected full-state geometry.  We query OSRM/OpenStreetMap for source-to-market road-network driving times from those constructed hub points to purchase municipalities \citep{openstreetmap2026,luxen2011osrm}.  The relevant measure is therefore not a straight-line distance but the shortest-route driving time implied by the OSM road network:
  \[
    T_m = \min_{h \in \mathcal{H}} \text{time}_{h,m},
  \]
  where $m$ indexes the purchase municipality, $h$ indexes the five hub states, and $\text{time}_{h,m}$ is measured in driving minutes.  The direction is hub-to-market because the instrument is meant to proxy supply-chain transportation costs; in a directed road network, this need not equal the reverse trip.  Operationally, these are continuous origin--destination driving times: the pairwise analogue of isochrone travel-time bands, but without coarsening municipalities into broad polygons.  The instrument is $\ln T_m$.  The identifying assumption is that longer road-network driving times proxy for transportation and supply-chain costs that shift local retail prices, without directly affecting the quantity a given buyer demands conditional on controls and fixed effects.

  \item \textbf{Lagged seizure-gravity indexes.}  Using monthly municipality-level marijuana-seizure data compiled by \citet{mucd2019seizures} from official federal responses to public-information requests, we construct gravity-weighted supply-disruption measures:
  \[
    Z_{m,t-1} = \sum_j \frac{\text{kg}_{j,t-1}}{1+\tau_{j,m}},
  \]
  where $\text{kg}_{j,t-1}$ is the kilograms of marijuana seized in municipality $j$ during the month before the purchase, and $\tau_{j,m}$ is the OSRM/OpenStreetMap driving time in minutes from seizure municipality $j$ to purchase municipality $m$.  The one-month lag allows supply disruptions to propagate through local retail markets rather than assuming a same-month pass-through.  The $1+\tau_{j,m}$ denominator avoids division by zero and imposes a simple inverse-driving-time decay: seizures that are closer in the source-to-market road network receive more weight than seizures that are farther away in driving time.  Same-municipality seizures are retained; when $j=m$, the driving time is zero and the denominator equals one, so local seizures enter with full weight rather than being discarded.  We use $\ln(1+Z_{m,t-1})$ as the instrument.  The index is merged to each transaction by municipality and purchase month.  Because purchases in the analytic sample run from January through September~2019, the one-month lag uses MUCD marijuana-seizure records from December~2018 through August~2019 and can be merged to all 3{,}293 transactions.  Its identifying logic is that nearby or local interdiction activity can disrupt supply chains and shift retail prices, while being plausibly unrelated to the idiosyncratic quantity chosen by an individual buyer after conditioning on controls and fixed effects.

  We report two variants.  The hub-only version restricts $j$ to seizure municipalities in the five production-hub states used in the OSM driving-time instrument, so it asks whether enforcement shocks in upstream production areas predict retail prices across destination markets.  The stricter outside-region version excludes seizure municipalities in the buyer's own macroregion, reducing exposure to local demand shocks at the cost of weaker first-stage power.  Both variants use only marijuana seizures, not crop destruction or seed confiscations.

  \item \textbf{External Hausman-type instruments.}  Following the logic of Hausman-style price instruments used in differentiated-product demand settings \citep{hausman1994competitive,nevo2001cereal}, we compute mean log prices in other markets within the same quality--month cell.  The preferred version excludes all transactions from the buyer's own state; a stricter version excludes the buyer's entire macroregion.  This construction avoids using prices from the same local state market as the endogenous regressor, while preserving common supply-cost variation across markets for the same quality tier and month.  The corresponding IV columns also include quality-by-month fixed effects, so identification comes from cross-market price covariation net of national quality-specific monthly shocks.
\end{enumerate}

For each IV strategy, we report first-stage F-statistics and assess instrument strength.  We complement the IV analysis with Anderson--Rubin confidence sets that are robust to weak instruments.  Wu--Hausman p-values are reported only when the first stage clears the conventional $F \geq 10$ threshold, because the test is not informative under weak identification.

\paragraph{Robustness.}
We subject the baseline estimates to several additional checks, keeping the state-by-month fixed effects and demographic controls wherever the specification permits direct comparability: (i)~trimming at the 1st and 99th percentiles of price and quantity to reduce the influence of outliers; (ii)~excluding round-number quantities to address heaping-induced measurement error; (iii)~adding municipality fixed effects to absorb time-invariant local heterogeneity; (iv)~using wild cluster bootstrap inference to account for the limited number of state clusters; and (v)~the \citet{oster2019unobservable} coefficient-stability test, which quantifies the degree of selection on unobservables required to explain away the estimated elasticity.


\section{Results}\label{sec:results}
% ============================================================
%  sections/results.tex  –  Results
% ============================================================

\subsection{Baseline OLS estimates}

Table~\ref{tab:main-ols} reports the core log--log estimates linking transaction quantity to unit price.  The pooled specification implies an elasticity of $-0.551$, and the coefficient becomes slightly more negative as controls are added.  Once product quality and demographics are included, the estimate is $-0.633$.  We use the state-by-month fixed-effects specification in column~(6) as the main benchmark because it keeps the full analytic sample while absorbing state-specific monthly shocks, the richest time-varying market controls available in the main OLS block.  In that benchmark, the elasticity is $-0.652$.  The preceding state-only and state-plus-month specifications produce very similar elasticities of $-0.643$, so the estimate remains in a narrow band as the specification absorbs richer geographic and temporal heterogeneity.  This stability suggests that the baseline association is not driven solely by broad local or seasonal shocks.

The magnitude is also consistent with the transaction-based cannabis literature.  It is close to the crowdsourced estimate of \citet{davis2016crowd} for the United States and falls comfortably within the range of geography-based estimates summarized in Section~\ref{sec:litreview}, including \citet{halcoussis2017geo}.  In substantive terms, the benchmark coefficient implies that a 10\% higher unit price is associated with roughly a 6.5\% smaller transaction quantity. 

\begin{table}[H]
\centering
\caption{OLS estimates of the price elasticity of demand for marijuana}
\label{tab:main-ols}
\centering
\resizebox{\linewidth}{!}{%
\begin{tabular}[t]{lcccccc}
\toprule
  & (1) Pooled & (2) +Quality & (3) +Demog. & (4) +State FE & (5) +Month FE & (6) St×Mo FE\\
\midrule
ln(Price) & -0.551*** & -0.615*** & -0.633*** & -0.643*** & -0.643*** & -0.652***\\
 & (0.016) & (0.017) & (0.017) & (0.020) & (0.021) & (0.022)\\
Good quality &  & 0.651*** & 0.617*** & 0.600*** & 0.595*** & 0.603***\\
 &  & (0.044) & (0.044) & (0.060) & (0.059) & (0.059)\\
Age &  &  & 0.026*** & 0.024*** & 0.024*** & 0.022***\\
 &  &  & (0.003) & (0.004) & (0.004) & (0.005)\\
Gender: Other &  &  & 0.077 & 0.121 & 0.138 & 0.165\\
 &  &  & (0.171) & (0.158) & (0.159) & (0.163)\\
Gender: Women &  &  & 0.016 & 0.027 & 0.016 & 0.034\\
 &  &  & (0.052) & (0.054) & (0.053) & (0.059)\\
Education: Middle school &  &  & -0.002 & 0.043 & 0.022 & 0.010\\
 &  &  & (0.236) & (0.220) & (0.229) & (0.206)\\
Education: High school &  &  & 0.025 & 0.084 & 0.063 & 0.070\\
 &  &  & (0.217) & (0.201) & (0.210) & (0.180)\\
Education: Bachelor &  &  & 0.287 & 0.347 & 0.322 & 0.348*\\
 &  &  & (0.217) & (0.208) & (0.218) & (0.184)\\
Education: Graduate &  &  & 0.339 & 0.415* & 0.397 & 0.415*\\
 &  &  & (0.226) & (0.234) & (0.242) & (0.209)\\
\midrule
Num.Obs. & 3293 & 3293 & 3293 & 3293 & 3293 & 3293\\
R2 & 0.288 & 0.326 & 0.367 & 0.402 & 0.408 & 0.448\\
R2 Adj. & 0.288 & 0.326 & 0.365 & 0.395 & 0.399 & 0.401\\
FE: state\_code &  &  &  & X & X & \\
FE: month\_str &  &  &  &  & X & \\
FE: state\_month &  &  &  &  &  & X\\
\bottomrule
\multicolumn{7}{l}{\rule{0pt}{1em}* p $<$ 0.1, ** p $<$ 0.05, *** p $<$ 0.01}\\
\multicolumn{7}{l}{\rule{0pt}{1em}Dependent variable: ln(quantity in grams).}\\
\multicolumn{7}{l}{\rule{0pt}{1em}Robust SEs in parentheses (cols 1--3); clustered at state level (cols 4--6).}\\
\end{tabular}
}
\end{table}


\subsection{Quality heterogeneity}

Table~\ref{tab:quality-segmented} shows that the average elasticity masks substantial heterogeneity across product tiers.  When the sample is split by perceived quality, seedless marijuana has an estimated elasticity of $-0.613$, whereas seeded marijuana has an elasticity of $-0.891$.  The pooled interaction specification delivers the same message.  In that column, the coefficient on $\ln(\text{Price})$ is $-0.873$ for the reference category, and the interaction term of $0.259$ implies an elasticity of about $-0.614$ for good-quality transactions.  The close match between the split-sample and interaction estimates suggests that the quality gradient is not an artifact of changing samples.  All three columns retain the same state-by-month fixed effects and demographic controls as the benchmark, so the comparison isolates quality-tier heterogeneity rather than changes in the control set.

One interpretation is that seeded marijuana functions as a less differentiated lower-tier product, for which consumers may have closer substitutes across sellers.  Seedless marijuana, by contrast, may attract buyers with stronger preferences for quality or fewer close substitutes within the same tier.  Although the descriptive statistics show a large price gap between quality tiers, the slope estimates still point to greater price responsiveness in the lower-quality segment.

\begin{table}[H]
\centering
\caption{Quality-segmented demand estimates}
\label{tab:quality-segmented}
\centering
\begin{tabular}[t]{lccc}
\toprule
  & (7a) Good only & (7b) Bad only & (8) Interaction\\
\midrule
ln(Price) & -0.613*** & -0.891*** & -0.873***\\
 & (0.022) & (0.041) & (0.036)\\
Good quality &  &  & -0.019\\
 &  &  & (0.118)\\
ln(Price) x Good &  &  & 0.259***\\
 &  &  & (0.037)\\
Age & 0.023*** & 0.022*** & 0.023***\\
 & (0.005) & (0.007) & (0.005)\\
Gender: Other & 0.174 & 0.305 & 0.161\\
 & (0.208) & (0.184) & (0.163)\\
Gender: Women & -0.005 & 0.171 & 0.037\\
 & (0.059) & (0.113) & (0.059)\\
Education: Middle school & -0.124 & 0.193 & 0.017\\
 & (0.262) & (0.242) & (0.220)\\
Education: High school & -0.021 & 0.234* & 0.079\\
 & (0.270) & (0.116) & (0.196)\\
Education: Bachelor & 0.270 & 0.496*** & 0.353*\\
 & (0.263) & (0.057) & (0.198)\\
Education: Graduate & 0.345 & 0.539*** & 0.418*\\
 & (0.294) & (0.135) & (0.219)\\
\midrule
Num.Obs. & 2636 & 657 & 3293\\
R2 & 0.442 & 0.644 & 0.454\\
R2 Adj. & 0.385 & 0.514 & 0.408\\
\bottomrule
\multicolumn{4}{l}{\rule{0pt}{1em}* p $<$ 0.1, ** p $<$ 0.05, *** p $<$ 0.01}\\
\multicolumn{4}{l}{\rule{0pt}{1em}Dependent variable: ln(quantity in grams).}\\
\multicolumn{4}{l}{\rule{0pt}{1em}All columns include state-by-month FE and demographics. SEs clustered at state level.}\\
\end{tabular}
\end{table}


\subsection{Instrumental-variables diagnostics}

Table~\ref{tab:iv-combined} summarizes the IV diagnostics using the OSM/OSRM driving-time matrices and MUCD marijuana-seizure records described in Section~\ref{sec:data}.  The table follows the same order as the empirical strategy: OSM driving time from production hubs, lagged seizure-gravity indexes, and external Hausman-type instruments.  All columns include the benchmark state-by-month fixed effects and demographic controls; the external Hausman columns also include quality-by-month fixed effects, so the standalone quality indicator is absorbed in those specifications.  The first geography-based diagnostic replicates the production-hub distance logic used in prior cannabis studies with OSM driving time rather than straight-line distance.  It remains uninformative in this setting: the IV coefficient is $0.793$ with a standard error of $2.647$.  The one-month-lag hub-only marijuana-seizure gravity instrument yields $-0.653$, nearly identical to the benchmark OLS estimate, while the stricter outside-region seizure-gravity instrument yields $-0.680$ but with a larger standard error.  The two external Hausman instruments, which exclude the buyer's own state or macroregion when averaging prices within quality--month cells, yield negative elasticities of $-0.712$ and $-1.025$, respectively.  Because purchases in the estimation sample run from January through September~2019, the one-month-lag seizure-gravity specifications use MUCD marijuana seizures from December~2018 through August~2019 and can be merged to all 3{,}293 transactions.

\begin{table}[H]
\centering
\caption{OLS and IV estimates: alternative identification strategies}
\label{tab:iv-combined}
\centering
\resizebox{\linewidth}{!}{%
\begin{tabular}[t]{lcccc}
\toprule
  & (1) OLS & (2) IV: OSM driving time & (3) IV: Hausman & (4) IV: gravity\\
\midrule
ln(Price) & -0.652*** &  &  & \\
 & (0.022) &  &  & \\
ln(Price) [IV] &  & 0.793 & -0.648*** & -0.735***\\
 &  & (2.647) & (0.068) & (0.248)\\
Good quality & 0.603*** & -0.688 & 0.590*** & 0.677**\\
 & (0.059) & (2.363) & (0.069) & (0.250)\\
Age & 0.022*** & 0.025*** & 0.022*** & 0.022***\\
 & (0.005) & (0.009) & (0.005) & (0.005)\\
Gender: Other & 0.165 & 0.235 & 0.213 & 0.161\\
 & (0.163) & (0.345) & (0.178) & (0.165)\\
Gender: Women & 0.034 & 0.024 & 0.029 & 0.035\\
 & (0.059) & (0.135) & (0.059) & (0.057)\\
Education: Middle school & 0.010 & 0.485 & -0.013 & -0.018\\
 & (0.206) & (0.942) & (0.207) & (0.214)\\
Education: High school & 0.070 & 0.358 & 0.064 & 0.053\\
 & (0.180) & (0.587) & (0.180) & (0.177)\\
Education: Bachelor & 0.348* & 0.372 & 0.354* & 0.347*\\
 & (0.184) & (0.291) & (0.183) & (0.180)\\
Education: Graduate & 0.415* & 0.195 & 0.429** & 0.427*\\
 & (0.209) & (0.562) & (0.208) & (0.219)\\
\midrule
Num.Obs. & 3293 & 3293 & 3173 & 3293\\
R2 & 0.448 & 0.136 & 0.133 & 0.137\\
R2 Adj. & 0.401 & 0.063 & 0.075 & 0.064\\
\bottomrule
\multicolumn{5}{l}{\rule{0pt}{1em}* p $<$ 0.1, ** p $<$ 0.05, *** p $<$ 0.01}\\
\multicolumn{5}{l}{\rule{0pt}{1em}Dependent variable: ln(quantity in grams).}\\
\multicolumn{5}{l}{\rule{0pt}{1em}All columns include state-by-month FE and demographic controls.}\\
\multicolumn{5}{l}{\rule{0pt}{1em}OSM driving-time IV: instrument = ln(OSM driving time from nearest constructed production-hub point).}\\
\multicolumn{5}{l}{\rule{0pt}{1em}Hausman IV: instrument = leave-one-out mean ln(price) within state-quality-month cell.}\\
\multicolumn{5}{l}{\rule{0pt}{1em}Gravity IV: instrument = ln(1 + one-month-lag seizure-gravity index).}\\
\multicolumn{5}{l}{\rule{0pt}{1em}For weak geography-based IVs, inference should rely on Anderson-Rubin confidence sets rather than conventional significance stars.}\\
\multicolumn{5}{l}{\rule{0pt}{1em}SEs clustered at state level.}\\
\multicolumn{5}{l}{\rule{0pt}{1em}* p $<$ 0.1, ** p $<$ 0.05, *** p $<$ 0.01.}\\
\end{tabular}
}
\end{table}


Table~\ref{tab:iv-diagnostics} reports the corresponding first-stage and weak-instrument diagnostics, while repeating the IV coefficient from Table~\ref{tab:iv-combined} for readability.  The OSM production-hub instrument has a first-stage F-statistic of only $0.60$, well below conventional thresholds for reliable IV inference \citep{staiger1997instrumental,stock2005testing}.

The seizure-gravity diagnostics show a useful strength--exclusion tradeoff.  The hub-only lagged marijuana-seizure gravity instrument clears the conventional first-stage threshold ($F=15.64$), produces an IV estimate of $-0.653$, and has an Anderson--Rubin confidence set of $[-1.08,-0.41]$.  This is the closest IV estimate to the benchmark and is consistent with the interpretation that upstream enforcement shocks in production hubs shift retail prices.  However, the hub-only design still allows enforcement variation from broad production regions that may be connected to other market conditions.  The stricter outside-region seizure instrument addresses that concern more directly, but its first stage falls to $F=6.82$ and the Anderson--Rubin confidence set spans the full grid reported in the table.  For weak specifications, inference should be read through the Anderson--Rubin confidence sets rather than conventional second-stage significance stars \citep{anderson1949estimation}.

The external Hausman instruments are statistically stronger: excluding the buyer's own state gives $F=51.61$, and excluding the buyer's macroregion gives $F=21.73$.  The associated Anderson--Rubin intervals exclude zero, but the exclusion restriction remains an assumption about common cross-market price shocks rather than an empirical fact.  The Wu--Hausman p-values of $0.566$ and $0.152$ therefore indicate that OLS and these IV estimates are not statistically distinguishable under those maintained assumptions; they do not prove exogeneity.

\begin{table}[H]
\centering
\caption{Instrumental-variables estimates and diagnostics}
\label{tab:iv-diagnostics}
\centering
\resizebox{\linewidth}{!}{%
\begin{tabular}[t]{lccc}
\toprule
Diagnostic & OSM driving-time IV & Hausman IV & Seizure-gravity IV\\
\midrule
First-stage F & 0.60 & 169.09 & 7.07\\
IV estimate & 0.793 (2.647) & -0.648 (0.068) & -0.735 (0.248)\\
Wu--Hausman p-value & 0.335 & 0.951 & 0.740\\
Anderson--Rubin 95\% CI & [-2.50, 0.50] & [-0.79, -0.52] & [-2.50, 0.42]\\
\bottomrule
\multicolumn{4}{l}{\rule{0pt}{1em}All diagnostics use the state-by-month FE specifications in Table~\ref{tab:iv-combined}.}\\
\multicolumn{4}{l}{\rule{0pt}{1em}IV estimate entries reproduce the coefficient on ln(Price) [IV] and clustered SE from Table~\ref{tab:iv-combined}.}\\
\multicolumn{4}{l}{\rule{0pt}{1em}Seizure-gravity diagnostics use one-month-lagged MUCD marijuana seizures from December 2018--August 2019.}\\
\multicolumn{4}{l}{\rule{0pt}{1em}For weak geography-based IVs, inference should rely on Anderson-Rubin confidence sets rather than conventional significance stars.}\\
\multicolumn{4}{l}{\rule{0pt}{1em}OSM and seizure-gravity have first-stage F-statistics below 10; Hausman has a strong first stage but relies on a leave-one-out exclusion restriction.}\\
\end{tabular}
}
\end{table}


Taken together, the IV evidence supports a cautious reading of the benchmark.  The weak OSM result shows that production-hub driving time alone is not enough to identify a strong first stage in this sample.  The hub-only lagged seizure-gravity instrument reproduces the OLS magnitude with a first stage above 10, while the stricter outside-region version shows that stronger exclusion restrictions can come with weaker first-stage power.  The external Hausman instruments show that cross-market price variation produces negative elasticities of similar order.  We therefore treat the IV block as triangulation around the transaction-level price--quantity gradient, not as a standalone causal proof that supersedes the benchmark OLS specification.

\subsection{Robustness and sensitivity}

The baseline elasticity remains negative and economically meaningful under several robustness checks.  Table~\ref{tab:robustness} first shows that trimming the 1st and 99th percentiles of price and quantity moves the benchmark estimate only from $-0.652$ to $-0.620$.  Excluding round-number quantities, a direct response to heaping documented in the raw quantity distribution, yields an elasticity of $-0.576$.  Both checks therefore reduce the magnitude slightly, but neither eliminates the central negative price--quantity gradient.

\begin{table}[H]
\centering
\caption{Robustness: outlier trimming and heaping}
\label{tab:robustness}
\centering
\begin{tabular}[t]{lccc}
\toprule
  & (6) Benchmark & (9) Trimmed & (10) No heap\\
\midrule
ln(Price) & -0.652*** & -0.620*** & -0.576***\\
 & (0.022) & (0.022) & (0.028)\\
Good quality & 0.603*** & 0.586*** & 0.511***\\
 & (0.059) & (0.058) & (0.088)\\
Age & 0.022*** & 0.023*** & 0.009\\
 & (0.005) & (0.005) & (0.009)\\
Gender: Other & 0.165 & 0.071 & -0.250\\
 & (0.163) & (0.132) & (0.478)\\
Gender: Women & 0.034 & -0.012 & -0.002\\
 & (0.059) & (0.062) & (0.057)\\
Education: Middle school & 0.010 & -0.035 & -0.540***\\
 & (0.206) & (0.203) & (0.081)\\
Education: High school & 0.070 & 0.008 & -0.348***\\
 & (0.180) & (0.186) & (0.071)\\
Education: Bachelor & 0.348* & 0.271 & -0.170*\\
 & (0.184) & (0.195) & (0.087)\\
Education: Graduate & 0.415* & 0.313 & -0.002\\
 & (0.209) & (0.227) & (0.125)\\
\midrule
Num.Obs. & 3293 & 3198 & 1417\\
R2 & 0.448 & 0.406 & 0.401\\
R2 Adj. & 0.401 & 0.356 & 0.301\\
\bottomrule
\multicolumn{4}{l}{\rule{0pt}{1em}* p $<$ 0.1, ** p $<$ 0.05, *** p $<$ 0.01}\\
\multicolumn{4}{l}{\rule{0pt}{1em}Dependent variable: ln(quantity in grams).}\\
\multicolumn{4}{l}{\rule{0pt}{1em}All columns include state-by-month FE plus demographics. SEs clustered at state level.}\\
\multicolumn{4}{l}{\rule{0pt}{1em}Trimmed = drop 1st/99th percentile of price and quantity.}\\
\multicolumn{4}{l}{\rule{0pt}{1em}No heap = exclude round-number quantities (multiples of 5 or near 28g).}\\
\multicolumn{4}{l}{\rule{0pt}{1em}* p $<$ 0.1, ** p $<$ 0.05, *** p $<$ 0.01.}\\
\end{tabular}
\end{table}


Table~\ref{tab:muni-fe} provides an additional spatial credibility check by adding municipality fixed effects to the state-by-month benchmark.  This comparison is not purely a change in fixed effects, since the municipality-fixed-effects specification drops municipalities with one observation and uses $3{,}127$ observations rather than $3{,}293$.  Even in that restricted sample, the elasticity changes only from $-0.652$ to $-0.664$.  The augmented specification therefore absorbs both state-month shocks and time-invariant municipality heterogeneity, yet the estimated elasticity changes by only about one hundredth.

\begin{table}[H]
\centering
\caption{Demand estimates with alternative fixed-effects structures}
\label{tab:muni-fe}
\centering
\resizebox{\linewidth}{!}{%
\begin{tabular}[t]{lcc}
\toprule
  & (1) St×Mo FE & (2) St×Mo + Muni FE\\
\midrule
ln(Price) & -0.652*** & -0.664***\\
 & (0.022) & (0.021)\\
Good quality & 0.603*** & 0.571***\\
 & (0.059) & (0.053)\\
Age & 0.022*** & 0.022***\\
 & (0.005) & (0.004)\\
Gender: Other & 0.165 & 0.073\\
 & (0.163) & (0.143)\\
Gender: Women & 0.034 & -0.008\\
 & (0.059) & (0.060)\\
Education: Middle school & 0.010 & -0.215\\
 & (0.206) & (0.199)\\
Education: High school & 0.070 & -0.162\\
 & (0.180) & (0.194)\\
Education: Bachelor & 0.348* & 0.092\\
 & (0.184) & (0.198)\\
Education: Graduate & 0.415* & 0.175\\
 & (0.209) & (0.231)\\
\midrule
Num.Obs. & 3293 & 3127\\
R2 & 0.448 & 0.494\\
R2 Adj. & 0.401 & 0.413\\
\bottomrule
\multicolumn{3}{l}{\rule{0pt}{1em}* p $<$ 0.1, ** p $<$ 0.05, *** p $<$ 0.01}\\
\multicolumn{3}{l}{\rule{0pt}{1em}Dependent variable: ln(quantity in grams).}\\
\multicolumn{3}{l}{\rule{0pt}{1em}All columns include state-by-month FE and demographic controls.}\\
\multicolumn{3}{l}{\rule{0pt}{1em}Col 2 additionally includes municipality FE and restricts to municipalities with at least 2 observations.}\\
\multicolumn{3}{l}{\rule{0pt}{1em}SEs clustered at state level.}\\
\end{tabular}
}
\end{table}


With only 32 state clusters, we also complement conventional clustered standard errors with a wild cluster bootstrap in the sense of \citet{cameron2008bootstrap} and \citet{roodman2019fast}.  Table~\ref{tab:bootstrap-inference} reports the conventional state-clustered interval, the wild cluster bootstrap interval using Webb weights and 9{,}999 replications, and a pairs cluster bootstrap sensitivity check with the same number of replications.  The wild-bootstrap interval remains tightly centered around the benchmark elasticity and excludes zero by a wide margin, so the main inference does not appear to be an artifact of conventional clustered standard errors.

\begin{table}[H]
\centering
\caption{Bootstrap inference for the benchmark elasticity}
\label{tab:bootstrap-inference}
\centering
\begin{tabular}[t]{lccc}
\toprule
Method & Estimate & SE & 95\% CI\\
\midrule
State-clustered SE & -0.652 & 0.023 & [-0.696, -0.607]\\
Wild cluster bootstrap & -0.652 & -- & [-0.715, -0.602]\\
Pairs cluster bootstrap & -0.653 & 0.023 & [-0.704, -0.610]\\
\bottomrule
\multicolumn{4}{l}{\rule{0pt}{1em}Benchmark model: state-by-month FE, quality, age, gender, and education controls.}\\
\multicolumn{4}{l}{\rule{0pt}{1em}Wild bootstrap uses Webb weights, 9,999 replications, and state clusters.}\\
\multicolumn{4}{l}{\rule{0pt}{1em}Pairs cluster bootstrap resamples states with replacement across 1,000 replications.}\\
\end{tabular}
\end{table}


Finally, Table~\ref{tab:oster-bounds} reports the coefficient-stability exercise of \citet{oster2019unobservable}.  The calculation uses the short and long specifications from columns~(1) and~(6) of Table~\ref{tab:main-ols}, with $R_{\max}$ set to $1.3R^2_{\text{long}}$.  Because adding observables moves the elasticity away from zero, the raw $\delta^{*}$ needed to set the coefficient to zero is negative; the table reports its absolute value, $|\delta^{*}|=7.70$.  Thus, explaining away the benchmark would require unobserved selection that is large and works in the opposite direction from the observed controls.  This is not a proof of identification, but it makes an omitted-variable explanation of the benchmark elasticity less plausible.

\begin{table}[H]
\centering
\caption{Oster coefficient-stability diagnostics}
\label{tab:oster-bounds}
\centering
\begin{tabular}[t]{p{0.62\linewidth}r}
\toprule
Statistic & Estimate\\
\midrule
Short-specification elasticity & -0.551\\
Short-specification $R^2$ & 0.288\\
Long-specification elasticity & -0.652\\
Long-specification $R^2$ & 0.448\\
$R_{\max}=1.3R^2_{long}$ & 0.582\\
Bias-adjusted elasticity ($\delta=1$) & -0.736\\
$|\delta^*|$ to set elasticity to zero & 7.70\\
\bottomrule
\multicolumn{2}{p{0.78\linewidth}}{\footnotesize \rule{0pt}{1em}Short specification includes only ln(price); long specification adds controls plus state-by-month FE.}\\
\multicolumn{2}{p{0.78\linewidth}}{\footnotesize \rule{0pt}{1em}The table reports the absolute value of $\delta^*$; the raw calculation is negative because coefficients are negative.}\\
\end{tabular}
\end{table}



\section{Conclusions}\label{sec:conclusions}
% ============================================================
%  sections/conclusions.tex  –  Conclusions
% ============================================================

This paper provides the first transaction-level estimate of the price elasticity of marijuana demand in Mexico.  Using 3{,}293 crowdsourced purchases observed across all 32 states and 351 municipalities, the preferred state-by-month fixed-effects specification yields an elasticity of $-0.652$.  The estimate implies that a 10\% higher unit price is associated with a 6.5\% smaller quantity purchased in a completed transaction.  This is an intensive-margin estimate: it speaks to transaction size among observed buyers, not to participation, purchase frequency, total individual consumption, or aggregate market demand.  Within that margin, the result is economically meaningful but less than unit-elastic, close to transaction-based evidence from other cannabis markets \citep{davis2016crowd,halcoussis2017geo,riley2020southafrica} and within the broader range summarized by \citet{gallet2014monkey}.

A central finding is that product quality matters.  Seedless marijuana is substantially more expensive and less price responsive, with an estimated elasticity of about $-0.61$, while seeded marijuana is more elastic, at about $-0.89$.  This difference suggests that consumers do not treat all marijuana transactions as a homogeneous product.  For policy design, the quality gradient matters because weight-based, price-based, and potency-sensitive tax instruments may affect product tiers differently.  A uniform per-gram tax, for example, would represent a larger proportional price increase for lower-quality product, precisely the segment in which demand appears more price responsive.

The benchmark estimate is stable across the main sensitivity checks.  Trimming extreme price and quantity observations changes the elasticity only modestly, excluding visibly heaped quantities preserves a negative and significant estimate, and adding municipality fixed effects yields an elasticity of $-0.664$ in the restricted sample.  Wild cluster bootstrap inference also confirms that the main result is not an artifact of having only 32 state clusters.  The Oster coefficient-stability exercise reports $|\delta^{*}|=7.70$; because the raw delta is negative under the sign convention induced by the negative coefficients, this magnitude should be read as the strength of unobserved selection required to move the benchmark estimate to zero.  These exercises do not prove causal identification, but they make a purely spurious interpretation of the price--quantity relationship less plausible.

The instrumental-variable evidence should be read in the same spirit.  The OSM production-hub driving-time instrument remains too weak to be informative as a standalone source of identification.  The hub-only one-month-lag marijuana-seizure gravity instrument is the closest IV analogue to the benchmark: it clears the conventional first-stage threshold and yields an elasticity of $-0.653$, while a stricter outside-region seizure instrument weakens substantially.  External Hausman instruments that exclude the buyer's own state or macroregion have strong first stages and produce negative elasticities of similar order, but they still rely on the assumption that other-market price variation captures supply-side cost shocks rather than correlated demand shocks.  The IV block therefore triangulates the benchmark and helps bound plausible interpretations, but it is not presented as a definitive causal proof.

The policy implication is not that a particular Mexican cannabis tax rate follows mechanically from an elasticity of $-0.65$.  The estimate comes from illicit-market transactions before national recreational regulation, not from a legal market with observed pass-through, compliance costs, product standards, or substitution between legal and illegal suppliers.  Still, it is a necessary empirical input.  Less-than-unit elastic demand implies that price-based regulation can reduce transaction quantities, but the response is unlikely to be so large that the tax base disappears.  At the same time, the quality-tier heterogeneity cautions against designing cannabis taxes as if the market were a single homogeneous product.  Any regulatory framework should therefore combine elasticity evidence with explicit assumptions about pass-through, enforcement, public-health objectives, and the coexistence of legal and illegal supply.

Several limitations remain.  The survey is crowdsourced and was promoted online, so the sample overrepresents young, male, and more educated consumers relative to the full population of marijuana users.  Quantities and prices are self-reported, and unit prices are constructed from reported value and quantity, so measurement error and nonlinear package pricing cannot be ruled out even after trimming and heaping checks.  Future work could combine transaction-level purchase data with representative consumption surveys, regulated-market administrative records if legalization proceeds, or repeated observations of buyers before and after policy changes.  Until such data exist, the evidence here offers a transparent and empirically grounded estimate of marijuana demand in Mexico at the transaction margin most directly relevant for price, taxation, and market-design debates.


\clearpage
\appendix
\section{External Data Descriptive Statistics}\label{app:external-data}
% ============================================================
%  sections/appendix_external_data.tex  –  External Data Descriptive Statistics
% ============================================================

This appendix documents the auxiliary enforcement and driving-time inputs used in the IV diagnostics.  These sources do not enter the baseline OLS specifications; they are used only in the supply-shock instruments described in Section~\ref{sec:data}.

Table~\ref{tab:appendix-mucd} summarizes the MUCD marijuana-seizure observations used to construct the seizure-gravity instrument, not the full MUCD database.  MUCD assembled these data from official responses to public-information requests submitted to federal security and justice institutions, then harmonized the heterogeneous files before publishing the database \citep{mucd2019seizures}.  Because purchases in the estimation sample run from January through September~2019, the one-month-lag instrument uses MUCD seizure records from December~2018 through August~2019.  In that lag window, 479 municipality-month records enter the marijuana-seizure index, spanning 9 months, 30 states, and 188 municipalities.  The distribution is highly right-skewed: the median seizure is 11.0 kilograms, the 90th percentile is 326.2 kilograms, and the largest recorded seizure is 10{,}390.0 kilograms.  This skewness motivates using a gravity-weighted index rather than treating all seizure locations as equally informative.

\begin{table}[H]
\centering
\caption{MUCD marijuana-seizure records used in IV diagnostics}
\label{tab:appendix-mucd}
\begin{threeparttable}
\begin{tabular}{lr}
\toprule
Statistic & Value\\
\midrule
Seizure lag period used & December 2018--August 2019\\
Lag months used & 9\\
States with marijuana seizures & 30\\
Municipalities with marijuana seizures & 188\\
Municipality-month seizure records used & 479\\
Total marijuana seized in lag months (kg) & 104,714.0\\
Median seizure (kg) & 11.0\\
90th percentile seizure (kg) & 326.2\\
Maximum seizure (kg) & 10,390.0\\
\bottomrule
\end{tabular}
\begin{tablenotes}[flushleft]\small
\item \textit{Notes:} MUCD denotes M\'{e}xico Unido contra la Delincuencia. The table reports the municipality-month marijuana-seizure records in the one-month lag window used to construct the seizure-gravity instrument. Purchases in the estimation sample run from January through September 2019, so the lagged MUCD months used are December 2018 through August 2019. MUCD compiles official responses to public-information requests submitted to federal security and justice institutions.
\end{tablenotes}
\end{threeparttable}
\end{table}


Table~\ref{tab:appendix-osm-osrm} reports diagnostics for the OSRM driving-time matrices built from OSM road data.  The matrices are based on a road-network snapshot queried in April~2026, not a historical reconstruction of 2019 routing conditions.  The updated production-hub matrix links five constructed hub-state points, derived from INEGI state geometries, to the 351 purchase municipalities.  The seizure-gravity matrix links the 188 MUCD marijuana-seizure municipalities used in the lag window to the same purchase municipalities.  Driving times are one-way source-to-market driving times in minutes, so the instruments use persistent road-network connectivity rather than straight-line distances.

\begin{table}[H]
\centering
\small
\caption{OSM/OSRM driving-time matrix diagnostics}
\label{tab:appendix-osm-osrm}
\begin{threeparttable}
\resizebox{\linewidth}{!}{%
\begin{tabular}{lrrrrrrrr}
\toprule
Matrix & Origins & Dest. & Pairs & Mean & Median & P10 & P90 & Max\\
\midrule
Production-hub states & 5 & 351 & 1,755 & 881.9 & 830.1 & 318.7 & 1,467.7 & 3,466.7\\
Seizure municipalities used & 188 & 351 & 65,988 & 891.6 & 743.2 & 225.6 & 1,737.2 & 4,878.4\\
\bottomrule
\end{tabular}
}
\begin{tablenotes}[flushleft]\small
\item \textit{Notes:} Driving times are one-way source-to-market driving times in minutes computed with the Open Source Routing Machine (OSRM) on OpenStreetMap (OSM) road data. The production-hub matrix links five constructed hub-state points, derived from INEGI state geometries, to purchase municipalities; the seizure-gravity matrix reports MUCD marijuana-seizure municipalities in the one-month lag window linked to purchase municipalities. OSRM/OSM is treated as a road-network snapshot queried and cached in April 2026, not as a time series.
\end{tablenotes}
\end{threeparttable}
\end{table}



\printbibliography
\end{document}
